In this section we collect the general results on contour integration that are used to deform the
contours of integration appearing in the construction of the magic functions $a$ and $b$. There are
two kinds of results: deformation of \emph{rectangular} contours (used for integrands defined on
vertical strips), and homotopy invariance of curve integrals of \emph{closed one-forms} (used to
move integrals along segments through the M\"obius inversion $z \mapsto -1/z$).

\subsection{Deformation of rectangular contours}

The starting point is the classical statement that the integral of a holomorphic function over the
boundary of a rectangle vanishes.

\begin{lemma}\label{lemma:rect-boundary-zero}\lean{Complex.hzero}\leanok
  Let $x_1, x_2, y \in \R$ and let $f : \C \to E$ be continuous on the closed half-strip
  $[x_1, x_2] \times [y, \infty)$ and complex differentiable on its interior, off a countable set.
  Then for every $m \geq y$,
  \[
    \int_{x_1}^{x_2} f(x + iy) \,\dd x - \int_{x_1}^{x_2} f(x + im) \,\dd x
      + i\int_{y}^{m} f(x_2 + it) \,\dd t - i\int_{y}^{m} f(x_1 + it) \,\dd t = 0,
  \]
  i.e.\ the integral of $f$ over the boundary of the rectangle with corners $x_1 + iy$ and
  $x_2 + im$ vanishes.
\end{lemma}
\begin{proof}\leanok
  This is a reformulation of the Cauchy--Goursat theorem for rectangles, which is available in
  Mathlib as \texttt{Complex.integral\_boundary\_rect\_eq\_zero\_of\_differentiable\_on\_off\_countable}.
\end{proof}

Letting the height of the rectangle tend to infinity, we obtain a deformation result for
``open rectangles'': the integral over a horizontal edge can be traded for the two vertical
half-lines above its endpoints.

\begin{theorem}\label{thm:open-rect-deform}\uses{lemma:rect-boundary-zero}\lean{Complex.integral_boundary_open_rect_eq_zero_of_differentiable_on_off_countable_of_integrable_on}\leanok
  Let $x_1, x_2, y \in \R$ and let $f : \C \to E$ be continuous on
  $[x_1, x_2] \times [y, \infty)$ and complex differentiable on its interior, off a countable set.
  Assume that
  \begin{enumerate}
    \item $t \mapsto f(x_1 + it)$ and $t \mapsto f(x_2 + it)$ are integrable on $(y, \infty)$;
    \item $f(z) \to 0$ uniformly as $\Im z \to \infty$.
  \end{enumerate}
  Then
  \[
    \int_{x_1}^{x_2} f(x + iy) \,\dd x
      + i\int_{y}^{\infty} f(x_2 + it) \,\dd t - i\int_{y}^{\infty} f(x_1 + it) \,\dd t = 0.
  \]
\end{theorem}
\begin{proof}\leanok
  Apply Lemma~\ref{lemma:rect-boundary-zero} to the rectangle of height $m$ and let
  $m \to \infty$. The uniform decay of $f$ makes the top edge vanish in the limit, and the
  integrability hypotheses ensure that the truncated vertical integrals converge to the integrals
  over the half-lines.
\end{proof}

\begin{theorem}\label{thm:rect-deform-tendsto-top}\uses{lemma:rect-boundary-zero}\lean{Complex.rect_deform_of_tendsto_top}\leanok
  The conclusion of Theorem~\ref{thm:open-rect-deform} continues to hold if the uniform decay
  hypothesis is replaced by the weaker assumption that the top-edge integrals tend to zero:
  \[
    \lim_{m \to \infty} \int_{x_1}^{x_2} f(x + im) \,\dd x = 0.
  \]
\end{theorem}
\begin{proof}\leanok
  Same argument as Theorem~\ref{thm:open-rect-deform}, using the hypothesis directly to dispose of
  the top edge.
\end{proof}

\subsection{Curve integrals of scalar one-forms}

To deform contours that are not parallel to the axes, we work with curve integrals of one-forms.
Every scalar function gives rise to a ($\C$-linear) one-form, and integrating this one-form along a
path recovers the usual complex contour integral.

\begin{definition}\label{def:one-form}\lean{curveIntegral}\leanok
  A \emph{one-form} on $\C$ (with values in $\C$) is a map
  \[
    \omega : \C \to \C \to_{L[\C]} \C
  \]
  assigning to every point $z \in \C$ a continuous $\C$-linear map $\omega(z) : \C \to \C$. Given a
  path $\gamma : [0,1] \to \C$, the \emph{curve integral} of $\omega$ along $\gamma$ is
  \[
    \int_{\gamma} \omega := \int_0^1 \omega(\gamma(t))\bigl(\gamma'(t)\bigr) \,\dd t,
  \]
  whenever the derivative $\gamma'$ exists almost everywhere and the integrand is integrable.
\end{definition}

\begin{definition}\label{def:scalarOneForm}\uses{def:one-form}\lean{MagicFunction.scalarOneForm}\leanok
  For a function $F : \C \to \C$, the associated \emph{scalar one-form} is
  \[
    \omega_F(z) : v \mapsto v \cdot F(z), \qquad z, v \in \C,
  \]
  regarded as a map $\C \to \C \to_{L[\C]} \C$ into the continuous $\C$-linear endomorphisms of
  $\C$.
\end{definition}

\begin{definition}\label{def:ClosedOneFormOn}\uses{def:one-form}\lean{SpherePacking.Contour.ClosedOneFormOn}\leanok
  A one-form $\omega : \C \to \C \to_{L[\C]} \C$ is a \emph{closed one-form on} a set
  $s \subseteq \C$ if
  \begin{enumerate}
    \item $\omega$ is differentiable on $s$ and continuous on the closure of $s$
      (\texttt{DiffContOnCl}); and
    \item its derivative within $s$ is symmetric: for all $x \in s$ and all $u, v$ in the tangent
      cone of $s$ at $x$,
      \[
        \bigl(D_u \omega(x)\bigr)(v) = \bigl(D_v \omega(x)\bigr)(u).
      \]
  \end{enumerate}
\end{definition}

\begin{remark}\label{rem:closed-one-form-comparison}
  Definition~\ref{def:ClosedOneFormOn} is the standard notion of a closed differential one-form,
  stated in coordinate-free language. For a $C^1$ one-form $\omega$ on an \emph{open} set
  $s$, the exterior derivative is given by
  \[
    d\omega(x)(u, v) = \bigl(D_u \omega(x)\bigr)(v) - \bigl(D_v \omega(x)\bigr)(u),
  \]
  so the symmetry condition~(2) is exactly $d\omega = 0$; moreover, on an open set the tangent
  cone at every point is all of $\C$ and the within-derivative agrees with the full derivative, so
  the two definitions coincide. The formalized version differs in two technical respects: it makes
  sense on an arbitrary set $s$ (via within-derivatives and tangent cones), and it demands
  continuity of $\omega$ up to the closure $\overline{s}$ (\texttt{DiffContOnCl}), because the
  deformation homotopies used later touch the boundary of the wedge region (e.g.\ the corner
  $1 \in \partial W$). The bundling of these two hypotheses into a single structure is specific to
  this project; the Poincar\'e lemma in Mathlib
  (Theorem~\ref{thm:poincare-curveIntegral}) takes them as separate hypotheses.

  For the scalar one-forms of Definition~\ref{def:scalarOneForm} closedness reduces to
  holomorphy: for $\omega_F = F \, \dd z$ one has
  $\bigl(D_u \omega_F(z)\bigr)(v) = v \cdot D_u F(z)$, and the symmetry condition with $u = 1$,
  $v = i$ reads $D_i F(z) = i \cdot D_1 F(z)$, which is precisely the Cauchy--Riemann equations.
  Hence $\omega_F$ is closed on an open set $s$ if and only if $F$ is holomorphic on $s$ --- the
  classical fact underlying Cauchy's theorem, and the way the closedness hypothesis is discharged
  in applications.
\end{remark}

The fundamental deformation principle is the Poincar\'e lemma for curve integrals, available in
Mathlib: curve integrals of a closed one-form along homotopic paths agree, up to the integrals
along the two transversal boundary paths of the homotopy.

\begin{theorem}[Poincar\'e lemma for curve integrals]\label{thm:poincare-curveIntegral}\uses{def:one-form, def:ClosedOneFormOn}\lean{ContinuousMap.Homotopy.curveIntegral_add_curveIntegral_eq_of_diffContOnCl}\leanok
  Let $s \subseteq \C$, let $\gamma_1, \gamma_2 : [0,1] \to \C$ be paths, and let
  $\varphi : [0,1] \times [0,1] \to \C$ be a homotopy from $\gamma_1$ to $\gamma_2$ such that
  \begin{enumerate}
    \item $\varphi$ maps the open square $(0,1)^2$ into $s$;
    \item $\varphi$ is $C^2$ on the closed square; and
    \item $\omega$ is a closed one-form on $s$.
  \end{enumerate}
  Denote by $\sigma_0 := \varphi(0, \cdot)$ and $\sigma_1 := \varphi(1, \cdot)$ the transversal
  paths connecting the start points and the end points of $\gamma_1$ and $\gamma_2$ respectively.
  Then
  \[
    \int_{\gamma_1} \omega + \int_{\sigma_1} \omega
      = \int_{\gamma_2} \omega + \int_{\sigma_0} \omega.
  \]
  In particular, if the homotopy fixes both endpoints, the curve integrals of $\omega$ along
  $\gamma_1$ and $\gamma_2$ coincide.
\end{theorem}
\begin{proof}\leanok
  This is Mathlib's
  \texttt{ContinuousMap.Homotopy.curveIntegral\_add\_curveIntegral\_eq\_of\_diffContOnCl}.
\end{proof}

\begin{lemma}\label{lemma:curveIntegral-segment}\uses{def:scalarOneForm}\lean{SpherePacking.Integration.integral_dir_mul_muIoc01_eq_curveIntegral_segment}\leanok
  Let $F : \C \to \C$ and let $a, b \in \C$. Parametrizing the straight segment from $a$ to $b$ by
  $z(t) = (1-t)a + tb$ for $t \in (0,1]$, we have
  \[
    \int_0^1 (b - a)\, F(z(t)) \,\dd t = \int_{[a,b]} \omega_F,
  \]
  where the right-hand side is the curve integral of the one-form $\omega_F$ along the segment.
\end{lemma}
\begin{proof}\leanok
  Immediate from the definition of the curve integral along a segment.
\end{proof}

\begin{lemma}[Change of variables on a segment]\label{lemma:curveIntegral-change-of-variables}\uses{def:scalarOneForm, lemma:curveIntegral-segment}\lean{SpherePacking.Contour.curveIntegral_segment_eq_neg_curveIntegral_segment_map'_of}\leanok
  Let $a, b \in \C$, let $f : \C \to \C$ be continuous on the segment $[a,b]$ and differentiable
  along its interior (i.e.\ $t \mapsto f(z(t))$ is differentiable for $t \in (0,1)$, where
  $z(t) = (1-t)a + tb$), and let
  $\Psi_r, \Psi'_r : \C \to \C$ be families of functions indexed by $r \in \R$. Suppose that for
  all $t \in (0,1)$,
  \[
    (b - a)\, \Psi_r(z(t))
      = -\frac{\dd}{\dd t}\bigl(f(z(t))\bigr) \cdot \Psi'_r\bigl(f(z(t))\bigr).
  \]
  Then
  \[
    \int_{[a,b]} \omega_{\Psi_r} = -\int_{f \circ [a,b]} \omega_{\Psi'_r},
  \]
  where $f \circ [a,b]$ denotes the image path of the segment under $f$.

  (In the formalization only continuity of $f$ is assumed: Lean's derivative operator
  \texttt{deriv} is total, taking the junk value $0$ at points of non-differentiability, so the
  displayed identity is meaningful as stated and implicitly carries the differentiability
  requirement.)
\end{lemma}
\begin{proof}\leanok
  Both sides are interval integrals over $t \in [0,1]$; the hypothesis identifies the integrands
  pointwise almost everywhere.
\end{proof}

\begin{definition}\label{def:mobiusInv}\lean{SpherePacking.mobiusInv}\leanok
  The \emph{M\"obius inversion} is the map
  \[
    \sigma : \C \to \C, \qquad \sigma(z) := -\frac{1}{z}.
  \]
  It is holomorphic away from $0$, with derivative $\sigma'(z) = 1/z^2$, and maps the upper
  half-plane $\h$ to itself.
\end{definition}

\begin{lemma}\label{lemma:curveIntegral-mobiusInv}\uses{def:mobiusInv, lemma:curveIntegral-change-of-variables}\lean{SpherePacking.MobiusInv.curveIntegral_segment_neg_inv}\leanok
  Let $a, b \in \C$ be such that the open segment $(a, b)$ lies in the upper half-plane and
  $\sigma$ is continuous on $[a,b]$. If two families $\Psi_r, \Psi'_r$ satisfy
  \[
    \Psi_r(z) = -\sigma'(z)\, \Psi'_r(\sigma(z)) \qquad \text{for all } z \in \h,
  \]
  then
  \[
    \int_{[a,b]} \omega_{\Psi_r} = -\int_{\sigma \circ [a,b]} \omega_{\Psi'_r}.
  \]
  The analogous statement with $\Psi_r(z) = \sigma'(z)\, \Psi'_r(\sigma(z))$ and no sign change on
  the right-hand side also holds (\texttt{curveIntegral\_segment\_pos\_inv}).
\end{lemma}
\begin{proof}\leanok
  Apply Lemma~\ref{lemma:curveIntegral-change-of-variables} with $f = \sigma$, computing
  $\frac{\dd}{\dd t} \sigma(z(t)) = (b - a)\,\sigma'(z(t))$ by the chain rule.
\end{proof}

\subsection{Homotopy invariance in the wedge region}

The key geometric object is a wedge-shaped open subset of the upper half-plane in which our
contours can be deformed into one another.

\begin{definition}\label{def:wedgeSet}\lean{SpherePacking.wedgeSet}\leanok
  The \emph{wedge set} is
  \[
    W := \setof{z \in \C}{0 < \Im z, \; \Re z - 1 < \Im z, \; 1 - \Re z < \Im z}.
  \]
  It is an open, convex subset of the upper half-plane, and its closure meets the real axis only
  at the point $1$.
\end{definition}

\begin{lemma}\label{lemma:wedge-contour-deform}\uses{def:wedgeSet, def:ClosedOneFormOn, def:mobiusInv, thm:poincare-curveIntegral}\lean{SpherePacking.Contour.perm_J12_contour_h1, SpherePacking.Contour.perm_J12_contour_h2}\leanok
  Let $\Psi'_r : \C \to \C$ be a family of functions such that $\omega_{\Psi'_r}$ is a closed
  one-form on $W$ for every $r \in \R$. Then
  \begin{align}
    \int_{\sigma \circ [-1,\, -1+i]} \omega_{\Psi'_r}
      + \int_{[\sigma(-1+i),\, 1+i]} \omega_{\Psi'_r}
      &= \int_{[1,\, 1+i]} \omega_{\Psi'_r},
      \label{eqn:wedge-deform-h1} \\
    \int_{\sigma \circ [-1+i,\, i]} \omega_{\Psi'_r}
      &= \int_{[1+i,\, i]} \omega_{\Psi'_r}
      + \int_{[\sigma(-1+i),\, 1+i]} \omega_{\Psi'_r}.
      \label{eqn:wedge-deform-h2}
  \end{align}
\end{lemma}
\begin{proof}\leanok
  The image of the segment $[-1, -1+i]$ under $\sigma$ and the segment $[1, 1+i]$ are homotopic
  inside $W$: the affine homotopy between them stays in $W$ by convexity of $W$
  (Definition~\ref{def:wedgeSet}) together with the fact that $\sigma$ maps both segments into
  $W$. Since $\sigma(-1) = 1$, both paths also share the endpoint $1$. The Poincar\'e lemma for
  curve integrals (Theorem~\ref{thm:poincare-curveIntegral})
  then yields~\eqref{eqn:wedge-deform-h1}. The identity~\eqref{eqn:wedge-deform-h2} follows by
  the same argument applied to the segments $[-1+i, i]$ and $[1+i, i]$, using $\sigma(i) = i$.
\end{proof}

Combining the change of variables with the homotopy invariance gives the two contour permutation
identities used in the study of the functions $a$ and $b$: the integrals over the two left
segments of the contour (from $-1$ to $-1+i$ to $i$) can be replaced by integrals over the two
right segments (from $1$ to $1+i$ to $i$).

\begin{theorem}\label{thm:perm-contour-neg}\uses{lemma:curveIntegral-mobiusInv, lemma:wedge-contour-deform}\lean{SpherePacking.perm_J12_contour_mobiusInv_wedgeSet}\leanok
  Let $\Psi_r, \Psi'_r : \C \to \C$ be families of functions such that
  \[
    \Psi_r(z) = -\sigma'(z)\, \Psi'_r(\sigma(z)) \qquad \text{for all } z \in \h,
  \]
  and such that $\omega_{\Psi'_r}$ is a closed one-form on $W$ for every $r$. Then
  \[
    \int_{[-1,\, -1+i]} \omega_{\Psi_r} + \int_{[-1+i,\, i]} \omega_{\Psi_r}
      = -\left( \int_{[1,\, 1+i]} \omega_{\Psi'_r} + \int_{[1+i,\, i]} \omega_{\Psi'_r} \right).
  \]
\end{theorem}
\begin{proof}\leanok
  By Lemma~\ref{lemma:curveIntegral-mobiusInv}, the integrals of $\omega_{\Psi_r}$ over the left
  segments equal minus the integrals of $\omega_{\Psi'_r}$ over their images under $\sigma$.
  Adding the two identities of Lemma~\ref{lemma:wedge-contour-deform}, the auxiliary integral
  along $[\sigma(-1+i), 1+i]$ cancels, giving the claim.
\end{proof}

\begin{theorem}\label{thm:perm-contour-pos}\uses{lemma:curveIntegral-mobiusInv, lemma:wedge-contour-deform}\lean{SpherePacking.perm_I12_contour_mobiusInv_wedgeSet}\leanok
  Under the hypotheses of Theorem~\ref{thm:perm-contour-neg}, but with the transformation law
  \[
    \Psi_r(z) = \sigma'(z)\, \Psi'_r(\sigma(z)) \qquad \text{for all } z \in \h,
  \]
  we have
  \[
    \int_{[-1,\, -1+i]} \omega_{\Psi_r} + \int_{[-1+i,\, i]} \omega_{\Psi_r}
      = \int_{[1,\, 1+i]} \omega_{\Psi'_r} + \int_{[1+i,\, i]} \omega_{\Psi'_r}.
  \]
\end{theorem}
\begin{proof}\leanok
  Identical to the proof of Theorem~\ref{thm:perm-contour-neg}, using the positive-sign variant of
  Lemma~\ref{lemma:curveIntegral-mobiusInv}.
\end{proof}
